\section{Inferring the Drains a City Will Not Give You}
\label{sec:sinkfield}

The largest known physics gap in the twin is that no storm sewer is modelled at all. The model
over-floods precisely where working drains exist. Mumbai will not hand out a drainage map, so the
question is whether the drains can be inferred from what a satellite can see.

Two independent validations had already said the twin knows \emph{how much} water and not
\emph{where}. Fitting more constants was not going to help, because the 16 per-land-class
multipliers of Section~\ref{sec:dem} change depth and not location. So instead of more constants we
give the model the missing object.

\subsection{Method}
After each solver step, water leaves a cell at a fitted rate. This is a sewer and not infiltration,
so it never fills the soil column and never counts as recharge. The field is
$\mathrm{softplus}(\theta) \times \mathrm{max\ rate} \times \mathrm{built}$, which forces it
positive and restricts it to built land. It is fitted with Adam through the checkpointed simulator
on 8 training storms and held out on 3, giving 65\,536 parameters per tile instead of 16. There is
freedom to move water, priors to keep it physical, and held-out storms to keep it honest.

Three design decisions are worth stating, because each one was a failure first.

Soft-Dice alone cannot train this field. The smooth wet indicator saturates in deeply flooded
cells, so the gradient vanishes exactly where the model is most wrong. We added depth hinges whose
slope is constant regardless of depth. The first run did 60 iterations and learned nothing.

The loss terms must share units. Expressed as volume fractions of the storm's rain they are
commensurate. With arbitrary weights the fit either drains everything, giving six times more
outflow than truth in a synthetic recovery test, or nothing at all.

A minimum-outflow prior charges for every cubic metre drained. This is what makes the volume
interpretable, because the fit then picks the \emph{smallest} drainage that explains the observed
extent. Every outflow number below is therefore a lower bound and not an estimate.

\subsection{Result one: held-out extent improves by predicting less water}
Across four Mumbai tiles the mean held-out CSI goes from 0.0435 to 0.0479, a gain of 10\,\%, while
predicted wet cells fall by 10\,\%. Per tile the gains are $+8\,\%$, $+2\,\%$, $+3\,\%$ and
$+25\,\%$.

Those are small numbers and we would not report them alone. What makes them worth reporting is the
direction of the trade. Section~\ref{sec:twin} established that CSI can be bought simply by
predicting more water. This is the opposite. Every tile improves while getting drier, which is why
the wetness is quoted beside every CSI.

\subsection{Result two: it is the spatial field, not the extra freedom}
The obvious objection is that 65\,536 parameters will beat 16 parameters at anything. We tested it
with a ladder on one tile, holding the loss, the storms and the held-out dates fixed and changing
only the parameter count.

\begin{table}[!t]
\centering
\caption{The controls ladder on Mumbai-East. Only the parameter count changes. A single city-wide
drain rate does nothing, the 16-parameter physics calibration makes held-out extent worse, and only
the spatial field improves it.}
\label{tab:ladder_sink}
\begin{tabular}{lcccc}
\toprule
model & params & train & held-out & held-out wet cells \\
\midrule
textbook physics          & 0      & 0.0428 & 0.0442 & 19\,691 \\
one uniform drain rate    & 1      & 0.0426 & 0.0439 & 19\,433 \\
per-class Manning/infil.  & 16     & 0.0354 & 0.0402 & 18\,612 \\
\textbf{per-cell field}   & \textbf{65\,536} & \textbf{0.0488} & \textbf{0.0478} & \textbf{17\,454} \\
\bottomrule
\end{tabular}
\end{table}

A single city-wide drain rate does nothing. The 16-parameter physics calibration makes held-out
extent \emph{worse}, because it fits the training storms by turning every multiplier up about 1.8
times, which is what ``change depth, not location'' looks like when you let it try harder. Only the
spatial field improves held-out extent. The gain comes from spatial structure and not from
parameter count, which is the claim the ladder exists to test.

\subsection{Result three: the decisive test, which it fails}
We replayed the 83 quarantined crowdsourced depth reports through both arms with identical
machinery, changing only the drain field. Skill against the best baseline moves from $-0.918$ to
$-0.921$. Correlation moves from $+0.037$ to $+0.032$. Root mean square error is 0.356\,m in both
arms. Nothing moves, and the differences are in the fourth decimal, far below meaningful at
$n = 83$.

The interpretation is clean. Learning where water \emph{leaves} the surface is not the same as
learning where it \emph{goes}. The field is fitted on radar extent, which is a wet or dry mask. That
constrains where drainage must exist and says almost nothing about the relative depth of two wet
streets. So it improves the quantity it was trained on and leaves the ranking failure untouched.
Co-location remains open after three independent attacks, being physics calibration, resolution
analysis and now drainage inversion.

\subsection{Result four: what the city pours out, and an external anchor}
Integrating the drained volume over each storm answers a question the project could not previously
ask. Because the prior selects the smallest drainage consistent with the radar, these are lower
bounds.

On the complete six-tile city with design storms, inferred surface outflow is at least
3.48 million\,m$^3$ at 25\,mm, 4.63 at 50\,mm, 5.21 at 100\,mm and 5.41 at 200\,mm. Flooded land
over the same range goes from 70.7 to 241.5\,km$^2$.

\textbf{The outflow saturates.} Eight times the rain buys 56\,\% more drained water, and the last
doubling from 100 to 200\,mm buys 4\,\%. The field is rate limited. Once the streets are wet for
the whole storm, extra rainfall cannot leave any faster, so it ponds. Note how differently the two
columns behave. Drained volume flattens while flooded area keeps climbing. That is the mechanism.
Beyond the drainage ceiling, additional rain has nowhere to go but outward across the ground. This
is a concrete and falsifiable statement about Mumbai's drainage limit, at roughly 5.4 million cubic
metres per storm however hard it rains.

\textbf{We can now check it against a real municipal measurement.} The six municipal stormwater
pumping stations are rated at 258\,m$^3$/s combined, a theoretical ceiling of 22.3 million\,m$^3$
per day with every pump at full rate. More useful than the rating is a measurement. The Brihanmumbai
Municipal Corporation reported pumping 16\,451.55 million litres, which is 16.45 million\,m$^3$,
between 16 and 19 August 2025 through those six stations. That is about 4.7 million\,m$^3$ per day.

Our inferred ceiling of at least 5.41 million\,m$^3$ per storm sits just above that daily pumped
volume, which is the right side to sit on. Our number is total surface outflow, and Mumbai has 186
outfalls, most of them gravity fed to the sea at low tide. The municipal figure counts only what six
pump houses lifted. A total drainage estimate should exceed pump-only throughput, and it does, by a
modest factor rather than an implausible one.

This is a consistency check and not yet a validation, for two reasons we state plainly. Our number
is a lower bound by construction, so agreement in magnitude cannot be pushed further than ``not
contradicted''. And the time bases differ, because ours is per storm and theirs is per day over a
three and a half day wet spell. Closing it properly needs the per-event pumped volume for a specific
storm that we also model, and two dates in our mask set qualify.

Design capacity is the wrong quantity to compare against, incidentally. The city's drainage master
plan is specified as a rate over an area, at 25\,mm/h existing and 50\,mm/h designed, which converts
to roughly 10.9 million\,m$^3$ per hour over the 437\,km$^2$ of the municipal area and so exceeds
our per-storm figure by two orders of magnitude. That comparison says nothing, because only 28 of
58 project phases were complete as of 2019 and the rate is nominal. The pumped volume is the number
with teeth. Inferred capacities themselves are physically plausible for urban drainage, averaging
0.3 to 3.0\,mm/h over built land and peaking at 4 to 50\,mm/h.
